% Plantilla de Pandoc para los trabajos del Máster en Computación Cuántica (UNIR).
%
% Reproduce material/Plantilla_LATEX_FCC_MUCC_UNIR/plantilla.tex y carga su estilo
% estilo_unir-1.sty sin modificarlo (REGLAS.md, sección 8). Solo se añaden:
%   - las definiciones que necesita el LaTeX que genera Pandoc (listas, tablas,
%     figuras, citas de citeproc);
%   - tres ajustes técnicos que eliminan advertencias sin cambiar el diseño del modelo
%     (marcados con «ajuste técnico»).
\documentclass[11pt,a4paper,spanish]{book}
% Ajuste técnico: codificación T1. El modelo usa OT1, que no tiene versalita en negrita
% (babel la sustituye con una advertencia) y no separa bien las palabras con tilde.
\usepackage[T1]{fontenc}
\usepackage{estilo_unir-1}

% Ajuste técnico: la cabecera del modelo tiene dos líneas (\rhead{\theauthor\\}) y
% fancyhdr avisa de que 14pt no bastan; es la altura que fancyhdr aplicaría de todos modos.
\setlength{\headheight}{27.2pt}
% Ajuste técnico: la portada y el cuerpo empiezan ambos en la página 1; sin esto,
% hyperref crea dos destinos con el mismo nombre.
\hypersetup{hypertexnames=false}
% Ajuste técnico: la clase book iguala el final de todas las páginas (\flushbottom); con
% interlineado 1,5 y párrafos largos, LaTeX no siempre puede estirar la página y avisa
% (Underfull \vbox). Con \raggedbottom cada página termina en su última línea.
\raggedbottom

% --- Definiciones para el LaTeX de Pandoc -----------------------------------------
\providecommand{\tightlist}{%
  \setlength{\itemsep}{0pt}\setlength{\parskip}{0pt}}
\setlength{\emergencystretch}{3em}
% Citas y referencias generadas por citeproc (APA 7): lista con sangría francesa.
\NewDocumentCommand\citeproctext{}{}
\NewDocumentCommand\citeproc{mm}{%
  \begingroup\def\citeproctext{#2}\cite{#1}\endgroup}
\newlength{\cslhangindent}
\setlength{\cslhangindent}{1.27cm}
\newlength{\csllabelwidth}
\setlength{\csllabelwidth}{3em}
\newenvironment{CSLReferences}[2]
 {\begin{list}{}{%
  \setlength{\itemindent}{0pt}
  \setlength{\leftmargin}{0pt}
  \setlength{\parsep}{0pt}
  \ifodd #1
   \setlength{\leftmargin}{\cslhangindent}
   \setlength{\itemindent}{-1\cslhangindent}
  \fi
  \setlength{\itemsep}{#2\baselineskip}}%
  % APA 7: referencias alineadas a la izquierda (sin justificar); así las URL/DOI
  % largas no estiran la línea.
  \raggedright
  % El modelo carga natbib (author-year), cuyo \bibitem no acepta las entradas de
  % citeproc; dentro de esta lista cada entrada es un \item normal.
  \renewcommand{\bibitem}[2][]{\item}}
 {\end{list}}
\newcommand{\CSLBlock}[1]{\hfill\break\parbox[t]{\linewidth}{\strut\ignorespaces#1\strut}}
\newcommand{\CSLLeftMargin}[1]{\parbox[t]{\csllabelwidth}{\strut#1\strut}}
\newcommand{\CSLRightInline}[1]{\parbox[t]{\linewidth - \csllabelwidth}{\strut#1\strut}}
\newcommand{\CSLIndent}[1]{\hspace{\cslhangindent}#1}

%---------------------------
%título del trabajo y autor
%---------------------------
\title{Las simetrías en la física}
\author{Julián Calderón Almendros}
\date{DD de MES de AAAA}
\profesor{Rodrigo Gil-Merino y Rubio}

\begin{document}
\renewcommand{\listfigurename}{Índice de figuras}
\renewcommand{\listtablename}{Índice de tablas}
\renewcommand{\contentsname}{Índice de contenidos}
\renewcommand{\figurename}{Figura}
\renewcommand{\tablename}{Tabla}

\maketitle

\frontmatter
\tableofcontents

\chapter{Resumen}
Resumen en español del trabajo (máximo 150 palabras): objetivo o
propósito, metodología, resultados y conclusiones.

\vspace{\baselineskip}
{\bf Palabras clave:} palabra clave 1, palabra clave 2, palabra clave 3

\chapter{Abstract}
English version of the Resumen (maximum 150 words).

\vspace{\baselineskip}
{\bf Keywords:} keyword 1, keyword 2, keyword 3

\mainmatter
\chapter{Introducción}\label{introducciuxf3n}

\section{Motivación y
justificación}\label{motivaciuxf3n-y-justificaciuxf3n}

\section{Planteamiento del trabajo}\label{planteamiento-del-trabajo}

\section{Estructura del trabajo}\label{estructura-del-trabajo}

\chapter{Contexto y estado de la
cuestión}\label{contexto-y-estado-de-la-cuestiuxf3n}

\section{Contexto histórico}\label{contexto-histuxf3rico}

Entre finales de la década de 1860 y comienzos de la de 1870, se
desarrolló un movimiento matemático trascendental protagonizado por
\textbf{Felix Klein} (1849--1925) y \textbf{Sophus Lie} (1842--1899).
Ambos coincidieron en Berlín en el invierno de 1869 a 1870 y en París en
la primavera de 1870, y de esa colaboración salió un trabajo conjunto
sobre las curvas planas que quedan invariantes bajo un sistema de
transformaciones lineales (Klein y Lie, 1871). De ahí procede la idea de
que la geometría debía definirse a través de las propiedades y
magnitudes que permanecen invariantes bajo un determinado grupo de
transformaciones. Esta visión la formuló Klein en solitario en 1872, en
un folleto impreso con motivo de su incorporación como catedrático a la
Universidad de Erlangen (no fue un discurso), conocido como el Programa
de Erlangen (Klein, 1872/1893). El programa ordenó las distintas
geometrías de la época (euclídea, afín, proyectiva\ldots) mediante sus
grupos de transformaciones: cada geometría estudia lo que deja
invariante su grupo, y cuanto mayor es el grupo, menos propiedades se
conservan. Su influencia inmediata fue escasa; se difundió a partir de
1890, con sus reediciones y traducciones. Hay que decir que Klein y Lie,
a pesar de tener temperamentos muy diferentes, colaboraron estrechamente
durante su juventud, desarrollando ideas conjuntas sobre curvas y
transformaciones geométricas. Sin embargo, con los años sus caminos se
bifurcaron: mientras Klein se orientó hacia una visión global,
institucional y unificadora de las matemáticas, Lie se concentró en la
profunda especialización analítica que dio origen a la teoría de los
grupos continuos de transformaciones (hoy llamados grupos de Lie), con
la que buscaba hacer para las ecuaciones diferenciales lo que la teoría
de Galois había hecho para las ecuaciones algebraicas. La expuso con
Friedrich Engel en \emph{Theorie der Transformationsgruppen}, publicada
en tres volúmenes entre 1888 y 1893 (Lie y Engel, 1888--1893). La
relación entre ambos se rompió en 1892--1893 por disputas de prioridad:
en el prólogo del tercer volumen, Lie afirmó públicamente que no era
discípulo de Klein, lo que terminó por distanciar profesional y
personalmente a ambos genios (Guijarro Regalado, 2007).

Esta idea de invariancia llenaba el ambiente matemático, y en paralelo
se desarrolló la teoría algebraica de invariantes: el estudio de las
expresiones formadas con los coeficientes de una forma (un polinomio
homogéneo) que no cambian, o que solo quedan multiplicadas por una
potencia del determinante, al hacer un cambio lineal de variables. En
Gran Bretaña la iniciaron Arthur Cayley (1845/2009) y James Joseph
Sylvester en la década de 1840, y el irlandés George Salmon (1859) la
sistematizó en sus libros de texto. La llamada «trinidad invariante» no
incluye a Salmon: Sylvester dio ese nombre al grupo que formaba con
Cayley y el francés Charles Hermite (Parshall, 2024). En Alemania, la
escuela principal, paralela y coetánea a la británica, fue la de
\textbf{Alfred Clebsch} (1833--1872) y \textbf{Paul Gordan}
(1837--1912), que desarrollaron el método simbólico para calcular
invariantes (Clebsch, 1872). Gordan y \textbf{Max Noether} (1844--1921)
enseñaron en Erlangen, por lo que a veces se habla de una escuela de
invariantes de Erlangen, que no debe confundirse con el Programa de
Erlangen de Klein (Clebsch, por su parte, nunca estuvo en Erlangen:
trabajó en Karlsruhe, Giessen y Gotinga). El problema central de la
teoría era de finitud: demostrar que todos los invariantes de una forma
se obtienen como polinomios de un número finito de ellos. El enfoque
tradicional era un trabajo minucioso y descomunal de cálculo, en el que
cada caso concreto (formas de un grado y un número de variables dados)
debía estudiarse por separado, calculando explícitamente sus
invariantes; así demostró Gordan (1868) la finitud para las formas
binarias (de dos variables) de cualquier grado. Con un enfoque mucho más
novedoso, \textbf{David Hilbert} (1862--1943) resolvió el problema para
cualquier número de variables con su \textbf{Teorema de la Base} (1888,
publicado en 1890): todo ideal de un anillo de polinomios en un número
finito de variables sobre un cuerpo es finitamente generado (Hilbert,
1890). Pero su demostración era de existencia: probaba que hay una base
finita de invariantes sin dar un método para construirla, y por eso
Gordan llegó a calificarla de teología: «esto no es matemática, es
teología» (Almira Picazo y Sabina de Lis, 2007; Derksen y Kemper, 2002,
p. 49; Weyl, 1944). En 1893 Hilbert publicó una demostración más
constructiva (Hilbert, 1893). Gordan acabó cambiando de opinión: afirmó
haberse convencido de que «incluso la teología tiene sus méritos»
(Klein, 1926/1979, p. 311) y en 1899 publicó una versión simplificada de
la demostración de Hilbert (Gordan, 1899). El resultado no hizo
innecesario el cálculo explícito de invariantes, pero, al resolver con
carácter general el problema de la finitud, prácticamente cerró la
teoría de invariantes y le quitó durante un tiempo su interés a este
campo matemático.

La fusión definitiva entre esta potente álgebra de invariantes y la
geometría de las transformaciones no tardaría en revolucionar la física.
A principios del siglo XX, la física teórica se encontraba en plena
ebullición con el desarrollo de la Teoría de la Relatividad General de
Albert Einstein. Fue en ese escenario donde las simetrías continuas
---aquellas transformaciones suaves y fluidas descritas originalmente
por los grupos de Lie, como las rotaciones en el espacio o el avance del
tiempo--- revelaron su verdadero poder. El problema surgió en Gotinga:
David Hilbert, que en noviembre de 1915 había presentado una formulación
variacional de la gravitación, contemporánea de las ecuaciones de campo
de Einstein (Hilbert, 1915), se encontró con que en la nueva teoría la
energía no parecía conservarse de la manera local y matemática
tradicional (Hilbert sostenía que en ella no había teoremas de energía
«propios»). Para estudiar este problema, David Hilbert y Felix Klein
invitaron a Gotinga en 1915 a la mayor experta en álgebra de invariantes
de la época: Emmy Noether (1882--1935), hija de Max Noether y discípula
directa de Paul Gordan. Noether había hecho en 1907 su tesis doctoral
sobre invariantes calculados a mano, un trabajo ingente: un sistema
completo de 331 covariantes de la forma ternaria de cuarto grado
(Noether, 1908). Trabajó en el problema de la energía entre 1915 y 1918,
y el resultado fue el artículo \emph{Invariante Variationsprobleme}, que
Klein presentó el 26 de julio de 1918 en la Sociedad de Ciencias de
Gotinga (Noether, 1918a). El artículo contiene dos teoremas. El primero
es el que hoy se conoce como Teorema de Noether, una de las piedras
angulares de la física moderna: en todo sistema cuyas ecuaciones se
obtienen de un principio variacional, a cada simetría continua de la
acción (un grupo de Lie de transformaciones) le corresponde una ley de
conservación. Así, la invariancia bajo traslaciones en el tiempo implica
la conservación de la energía; bajo traslaciones en el espacio, la del
momento lineal, y bajo rotaciones, la del momento angular. El segundo
teorema trata las simetrías que dependen de funciones arbitrarias, como
los cambios generales de coordenadas de la relatividad general: en ese
caso no aparecen cantidades conservadas independientes, sino identidades
entre las ecuaciones del movimiento. Es este segundo teorema el que
explica la forma peculiar que toma la conservación de la energía en la
relatividad general, la que había señalado Hilbert. El trabajo de
Noether reunió así las dos líneas anteriores: los grupos continuos de
Lie y los métodos de la teoría de invariantes.

Einstein, que seguía estos trabajos desde Berlín, dejó constancia de su
aprecio por ellos en dos cartas (Einstein, 1998; Tollmien, 2018). El 24
de mayo de 1918, al recibir un artículo anterior de Noether sobre la
construcción de invariantes (Noether, 1918b), escribió a Hilbert: «me
impresiona que se puedan abarcar estas cosas desde un punto de vista tan
general» (Einstein, 1998, documento 548). El 27 de diciembre de 1918, al
recibir \emph{Invariante Variationsprobleme}, escribió a Klein que
volvía a sentir como una gran injusticia que a Noether se le negara la
\emph{venia legendi}, el permiso para enseñar en la universidad, y
propuso dar un paso enérgico ante el ministerio (Einstein, 1998,
documento 677). Noether obtuvo la habilitación en 1919.

\chapter{Objetivos}\label{objetivos}

\chapter{Desarrollo del trabajo}\label{desarrollo-del-trabajo}

\chapter{Conclusiones}\label{conclusiones}

\chapter*{Bibliografía}\label{bibliografuxeda}
\addcontentsline{toc}{chapter}{Bibliografía}

\protect\phantomsection\label{refs}
\begin{CSLReferences}{1}{0}
\bibitem[\citeproctext]{ref-almira2007}
Almira Picazo, J. M. y Sabina de Lis, J. C. (2007). \emph{Hilbert:
Matemático fundamental}. Nivola.

\bibitem[\citeproctext]{ref-cayley1845}
Cayley, A. (2009). On the theory of linear transformations. En \emph{The
collected mathematical papers of {Arthur Cayley}} (Vol. 1, pp. 80--94).
Cambridge University Press.
\url{https://doi.org/10.1017/CBO9780511703676.014} (Obra original
publicada en 1845)

\bibitem[\citeproctext]{ref-clebsch1872}
Clebsch, A. (1872). \emph{Theorie der binären algebraischen {Formen}}.
B. G. Teubner.

\bibitem[\citeproctext]{ref-derksen2002}
Derksen, H. y Kemper, G. (2002). \emph{Computational invariant theory}.
Springer. \url{https://doi.org/10.1007/978-3-662-04958-7}

\bibitem[\citeproctext]{ref-einstein1998}
Einstein, A. (1998). \emph{The collected papers of {Albert Einstein}:
Vol. 8. The {Berlin} years: Correspondence, 1914--1918} (R. Schulmann,
A. J. Kox, M. Janssen, y J. Illy, Eds.). Princeton University Press.

\bibitem[\citeproctext]{ref-gordan1868}
Gordan, P. (1868). Beweis, dass jede {Covariante} und {Invariante} einer
binären {Form} eine ganze {Function} mit numerischen {Coefficienten}
einer endlichen {Anzahl} solcher {Formen} ist. \emph{Journal für die
reine und angewandte Mathematik}, \emph{69}, 323--354.
\url{https://doi.org/10.1515/crll.1868.69.323}

\bibitem[\citeproctext]{ref-gordan1899}
Gordan, P. (1899). Neuer {Beweis} des {Hilbertschen} {Satzes} über
homogene {Funktionen}. \emph{Nachrichten von der Gesellschaft der
Wissenschaften zu Göttingen, Mathematisch-Physikalische Klasse},
\emph{1899}, 240--242. \url{http://eudml.org/doc/58452}

\bibitem[\citeproctext]{ref-guijarro2007}
Guijarro Regalado, J. L. (2007). \emph{Lie: Más allá de la geometría}.
Nivola.

\bibitem[\citeproctext]{ref-hilbert1890}
Hilbert, D. (1890). Ueber die {Theorie} der algebraischen {Formen}.
\emph{Mathematische Annalen}, \emph{36}(4), 473--534.
\url{https://doi.org/10.1007/BF01208503}

\bibitem[\citeproctext]{ref-hilbert1893}
Hilbert, D. (1893). Ueber die vollen {Invariantensysteme}.
\emph{Mathematische Annalen}, \emph{42}(3), 313--373.
\url{https://doi.org/10.1007/BF01444162}

\bibitem[\citeproctext]{ref-hilbert1915}
Hilbert, D. (1915). Die {Grundlagen} der {Physik} ({Erste}
{Mitteilung}). \emph{Nachrichten von der Gesellschaft der Wissenschaften
zu Göttingen, Mathematisch-Physikalische Klasse}, \emph{1915}, 395--408.
\url{http://eudml.org/doc/58946}

\bibitem[\citeproctext]{ref-klein1893}
Klein, F. (1893). Vergleichende {Betrachtungen} über neuere geometrische
{Forschungen}. \emph{Mathematische Annalen}, \emph{43}(1), 63--100.
\url{https://doi.org/10.1007/BF01446615} (Obra original publicada en
1872)

\bibitem[\citeproctext]{ref-klein1979}
Klein, F. (1979). \emph{Development of mathematics in the 19th century}
(M. Ackerman, Trad.). Math Sci Press. (Obra original publicada en 1926)

\bibitem[\citeproctext]{ref-kleinlie1871}
Klein, F. y Lie, S. (1871). Ueber diejenigen ebenen {Curven}, welche
durch ein geschlossenes {System} von einfach unendlich vielen
vertauschbaren linearen {Transformationen} in sich übergehen.
\emph{Mathematische Annalen}, \emph{4}(1), 50--84.
\url{https://doi.org/10.1007/BF01443297}

\bibitem[\citeproctext]{ref-lieengel1893}
Lie, S. y Engel, F. (1888--1893). \emph{Theorie der
{Transformationsgruppen}} (Vols. 1--3). B. G. Teubner.

\bibitem[\citeproctext]{ref-noether1908}
Noether, E. (1908). Über die {Bildung} des {Formensystems} der ternären
biquadratischen {Form}. \emph{Journal für die reine und angewandte
Mathematik}, \emph{134}, 23--90.
\url{https://doi.org/10.1515/crll.1908.134.23}

\bibitem[\citeproctext]{ref-noether1918}
Noether, E. (1918a). Invariante {Variationsprobleme}. \emph{Nachrichten
von der Gesellschaft der Wissenschaften zu Göttingen,
Mathematisch-Physikalische Klasse}, \emph{1918}, 235--257.
\url{http://eudml.org/doc/59024}

\bibitem[\citeproctext]{ref-noether1918a}
Noether, E. (1918b). Invarianten beliebiger {Differentialausdrücke}.
\emph{Nachrichten von der Gesellschaft der Wissenschaften zu Göttingen,
Mathematisch-Physikalische Klasse}, \emph{1918}, 37--44.
\url{http://eudml.org/doc/59011}

\bibitem[\citeproctext]{ref-parshall2024}
Parshall, K. H. (2024). A convergence of paths: {Cayley}, {Hermite},
{Sylvester}, and early invariant theory. \emph{Revue d'histoire des
mathématiques}, \emph{30}, 233--260.
\url{https://doi.org/10.24033/rhm.316}

\bibitem[\citeproctext]{ref-salmon1859}
Salmon, G. (1859). \emph{Lessons introductory to the modern higher
algebra}. Hodges, Smith {\&} Co.

\bibitem[\citeproctext]{ref-tollmien2018}
Tollmien, C. (2018). „{Invariantentheorie} ist jetzt hier {Trumpf}``.
\emph{Physik in unserer Zeit}, \emph{49}(4), 176--182.
\url{https://doi.org/10.1002/piuz.201801513}

\bibitem[\citeproctext]{ref-weyl1944}
Weyl, H. (1944). David {Hilbert} 1862--1943. \emph{Obituary Notices of
Fellows of the Royal Society}, \emph{4}(13), 547--553.
\url{https://doi.org/10.1098/rsbm.1944.0006}

\end{CSLReferences}

\appendix

\chapter{Apéndices}\label{apuxe9ndices}

\end{document}
